% Revision/source crosswalk: analyses/018-2026-09-08-results-materials/observations/O014-manuscript-polish.md
% Evidence: Analysis 018 observations O001/O002/O003/O004/O007/O011.
\section{Discussion and conclusion}
\label{sec:conclusion}

Pressure, thresholds and site placement need to be evaluated together.
In the matched 14M comparisons, adding pressure ranges from improving both
loss and sparsity to exchanging quality for additional zeros. Extending the
intervention to internal attention can increase model-wide sparsity even
when FFN sparsity falls. The high-threshold seven-site recipe with pressure
improves both axes over the four-site recipe at all three sizes, but the
zero-threshold ordering changes. These results favor evaluating each
addition against its matched recipe, rather than choosing interventions
from activation sparsity alone.

The interpretation is supported by a connected pattern across experiments.
Matched comparisons identify how pressure's added effect changes with the
threshold; distribution measurements reveal which activations change; and
product counts explain their model-wide consequence. The high-threshold
A7-OL1 advantage recurs at all three sizes, while the zero-threshold reversals
identify where that relationship stops holding. Agreement across these
comparisons makes the result more informative than a favorable checkpoint
or an isolated sparsity percentage.

Product accounting and kernel timing answer different questions. Counting
zero-operand products shows where an intervention affects computation;
the sparsity ceiling explains how architecture changes that opportunity.
The kernel study confirms that some of it can reduce full-model latency,
but also shows that skipping attention work can cost more than executing it.
Fusion and sparse-path ablations are therefore necessary to attribute a
measured speedup.

These conclusions concern MiniPile pretraining at a fixed token budget,
with a lower learning rate at 410M. Larger models test complete recipes;
the distribution diagnostics cover 14M and describe responses rather than
isolating a causal spillover mechanism. The kernel results use one RTX5090,
full-sequence inference, native SDPA and matched fusion controls. Transfer
to cached decoding, larger models and other hardware remains untested.
